% ===========================================================================
% 01-Introduction.tex  <- thesis/Sections/Introduction.tex
%
% SCOPE, per the skeleton and the brief: the problem, the object of study, the
% four questions, the win condition, and a map of the argument. v1's
% introduction also rehearsed most of the chapter-4 results; those are cut,
% because a result stated before the reader knows what would count as a good
% one is not an introduction to it. Every number kept below appears in v1 with
% the same value, and none is added. No \needsaudit flags: this chapter touches
% none of the three audited quantities.
%
% REBUILT. Three structural changes, all deliberate:
%
% (1) The chapter-level \begin{question} box is GONE. A framed panel standing
%     before the first sentence of the thesis asked the reader to parse an
%     instrument before they had the problem. Its three loads are redistributed
%     into the running argument: the retrieval test lands in sec:intro-question
%     where the metric is defined, the "bar is a lookup, not chance" and the
%     verdict land in sec:intro-bar where the lookup is built. The opening is
%     now two paragraphs of prose. This chapter has no opener box by design; do
%     not restore one.
%
% (2) The run-in \paragraph heads are down from fifteen to one. The chapter
%     read as an outline that had been filled in. Each section is now one or
%     two movements that carry the reader forward without a head to re-orient
%     them every eighty words.
%
% (3) sec:intro-map now comes FIRST of the three closing sections. The four
%     questions are the spine of the investigation and must be stated before
%     either instrument that judges them. sec:intro-question follows, so that
%     NIR arrives as the answer to the second half of question one ("which
%     metrics survive that control") rather than as a settled premise the
%     reader is asked to accept three pages before the chapter asks whether the
%     ruler can be trusted. sec:intro-bar stays last: the lookup is not one of
%     the four questions, it is the bar the re-encoded arm fails to clear, so
%     it closes the chapter with the win condition, the verdict and the map.
%     All three labels are unchanged. sec:intro-question and sec:intro-bar are
%     \cref'd only from inside this chapter; sec:intro-map is \cref'd from
%     04b-representation.tex:7,31, and moving it earlier only renumbers it.
%
% The four questions themselves are the research arc, each arising from the
% answer to the one before: broken ruler -> drug-blind on the repaired one ->
% why (data, representation, or readout) -> does re-encoding the target repair
% it. v1's fourth question was the lookup baseline; that is the judge of the
% fourth answer, not a fifth line of enquiry, and it has been re-sited.
%
% RESTORED after the cutting pass, in the robustness paragraph that closes
% sec:intro-bar: the list of what a backbone replication must preserve (split,
% targets, sentinels, token budgets, within-plate galleries, training-only
% lookups, validity and completion checks), verbatim from v1. That paragraph
% used to carry the run-in head "What survives the changes", which presupposed
% changes the chapter had never named; the head is gone and the four changes
% are now introduced before they are listed. Without the list the chapter said
% that scale is untested but no longer said what would make a scale test valid,
% which is
% the auditable half of the claim. The same cut had orphaned "It leaves an
% auditable ladder" -- the pronoun read as the untested backbone; v1's "This
% thesis does not supply that replication, but it leaves..." is back.
%
% The last paragraph carries the forward pointer to the reference apparatus,
% which main.tex moved to the back. It cites \nameref{sec:howtoread}, the label
% on the "Reference apparatus" chapter opener in 00-HowToRead.tex. \nameref and
% not \cref because that opener is \chapter*, i.e. uncounted. The two
% sentence-initial references in that paragraph are \Cref, not \cref: the
% run-in \paragraph head ends in a period, so "Chapter 2 places..." opens a
% sentence. v1 had "Section~\ref{...}" there and read "Section 2 places...";
% the port to cleveref lost the capital and set "chapter 2 places...".
%
% TERMINOLOGY, sec:intro-bar, the three decomposition figures: the shares are
% SQUARED EUCLIDEAN NORMS, not variances, and the prose now says so. Nothing is
% centred anywhere in the computation, so ||s||^2 is a sum of squares about the
% origin. The source of the old wording is a misnomer: inference.energy_shares
% is docstringed "Exact variance decomposition" while computing ||r||^2/||s||^2
% with no mean subtracted. Do not restore "energy" here on the strength of that
% docstring. See \glentry{$\lVert\cdot\rVert$} in 00-HowToRead.tex, which fixes
% L2 for the document and forecloses the variance reading.
%
% All three numbers trace to RESULTS_cluster/vardecomp_matched.json, main.scope:
% residual_energy_share 0.6211 -> 62.1%, generic_energy_share 0.3899 -> 39.0%,
% cross_energy_share -0.01106 -> -0.011; energy_shares_sum is 1.0 to 9 places,
% which is what "closing the arithmetic" means. The JSON keys keep the old name
% and are NOT renamed by this pass -- they are a data contract, and renaming
% them would strand build_thesis_assets.py:62,66, which reads those paths.
%
% One caveat the sentence does not carry, and deliberately: each share is a
% MEAN OVER CONDITIONS of a per-condition ratio (variance_decomposition.py:634,
% np.mean(rn**2 / s2)), not the decomposition of one pooled shift vector. "Its"
% therefore reads a little tighter than the estimator warrants. The glossary
% entry states the mean-over-conditions caveat once, for the whole document,
% which is the right place for it; do not expand this sentence to repeat it.
%
% One line-breaking hint, no words: the checkpoint name carries \allowbreak
% after "Pythia-" as well as after "C2S-Scale-". The single break point was
% not usable -- taking it left ~105pt to stretch, worse than tolerance -- so
% TeX preferred a 27.5pt overfull line. Breaking after "Pythia-" leaves ~7pt.
% ===========================================================================

\chapter{Introduction}
\label{sec:introduction}
\framedecl{ER}

\noindent
Give a model an untreated cell and the name of a drug, and ask it what that
drug will do to that cell. That is the task this thesis evaluates. It is worth
answering because the answer is otherwise obtainable only by running the
experiment.

The thesis takes one published cell-sentence checkpoint, fine-tunes it on a
hundred-million-cell drug atlas, and asks whether the predictions it produces
identify the drug that produced them. They do not. Most of what follows is the
apparatus that makes that a finding and not an anecdote: a metric calibrated
before it was trusted, controls a drug-blind predictor cannot pass, and a bar
that needs no model at all.

% ---------------------------------------------------------------------------
\section{The truth this task predicts is never observed}
\label{sec:intro-task}

A drug that enters a cell changes what that cell transcribes. Which genes move,
by how much, and in which direction is the readout the field treats as its most
direct account of what a drug does. That is a working premise, not an
established fact. A substantial share of the treated conditions in the atlas
used here are statistically inert or not empirically discriminable at the tested
pseudobulk and cell budget (\cref{sec:methods-data}), and cannot support this
evaluation.

Reading one response is routine; reading every combination that might matter is
not, because a drug behaves differently across doses and cellular backgrounds,
and even the largest screens sample a small corner of the space. A computed
response would make the untested combinations inspectable, which is what
prediction is for. The task is simple to state: given the transcriptome of an
untreated cell and the identity of a drug, predict the transcriptome that cell
would have shown had the drug been applied.

It is harder to grade than to state, because the ground truth is
counterfactual. Sequencing destroys the cell it reads, so a prediction can only
be checked against a \emph{different} cell that received the drug. Paired
training data must therefore be assembled by matching cells, evaluation is
framed at population level, and distinguishing a drug effect from ordinary
cell-to-cell variation requires explicit controls. The units, the frames and
the controls that follow all descend from that one fact.

% ---------------------------------------------------------------------------
\section{The atlas contributes breadth; the inferential unit decides what can be
claimed}
\label{sec:intro-atlas}

Tahoe-100M~\cite{tahoe100m} makes that evaluation possible at unusual scale,
with roughly one hundred million single-cell transcriptomes spanning about
$1{,}100$ small molecules across $50$ cancer cell lines. Each condition, one
drug at one dose in one cell line on one plate, is a population and not a
single summary.\seealso{\cref{sec:methods-data}: the atlas layout, and which
unit each interval is clustered on}

A population per condition admits two questions. The conditional one, what a
drug does in a particular cellular context, is what the task asks. The prior one
is whether the apparent response is there to be predicted at all, and how much
of it reproduces across disjoint halves of the same treated well. Standard
evaluations blur the two; much of this thesis keeps them apart.

Scale does not remove design limitations. Conditions are spread unevenly across
drugs, doses, cell lines and treated wells, and plate membership can change
which alternatives a rank-based score sees. Cells from the same well are not
independent biological replicates, and neither are drug rows sharing a cell
line. The analysis therefore builds pseudobulk truths from disjoint cell halves
and keeps the well as a resampling unit.

% ---------------------------------------------------------------------------
\section{A model can be fluent in cell sentences without reading the drug}
\label{sec:intro-model}

Cell2Sentence~\cite{c2s} writes a cell as its gene names in descending order of
expression, turning it into text an ordinary language model can
consume.\gloss{cell sentence}{A cell written as its expressed
genes, named in descending order of expression.} It discards expression
magnitudes and keeps ranks; in exchange it needs no specialised architecture.
C2S-Scale~\cite{c2sscale} extends it to models pretrained on more than $50$
million cells from $825$ single-cell datasets, fine-tuned across
classification, generation, captioning and perturbation tasks. Its published
perturbation work covers cytokine stimulation and, for small molecules, bulk
L1000 profiles and SciPlex3, evaluated primarily as distributions and not
condition by condition.

The checkpoint studied here is
\texttt{C2S-Scale-}\allowbreak\texttt{Pythia-}\allowbreak\texttt{1b-pt}, the
pretrained Pythia-1b model~\cite{pythia} from that family, with no drug-treated
cells in its declared pretraining corpus. Fine-tuning uses several hundred
thousand Tahoe cells, not the entire atlas. Every arm is cold-started from that
checkpoint with identical hyperparameters, so within an experiment the training
target or objective is the only variable. The prompt supplies the cell line, the
drug, the dose, a mechanism annotation and the untreated control cell's own
sentence, and what the model returns for that prompt is the prediction that gets
scored.

The training target separates the two principal arms, and with it the frame each
is scored in. The first arm predicts the treated cell's own sentence and is
scored in the expression frame (\Eframe), where a prediction is a cell profile
ranked against other drugs on the same plate. The later arm predicts the
drug-specific residual, meaning what survives once the drug-agnostic mean
response, the \emph{generic}, is subtracted from the treated-minus-control
shift. It is scored in the residual frame (\Rframe), where a prediction is a
drug-specific signature ranked against other drugs in the same cell line. Both
frames put chance in the same place, and neither is interchangeable with the
other. Which is in force is declared at the head of every section.

The appeal of the encoding is also the scientific risk. Housekeeping genes
dominate every ranked profile, treated cells retain most of the state of their
untreated lineage, and many perturbations share a generic response programme. A
model can therefore produce a plausible treated cell while remaining almost
indifferent to which drug was named.

\begin{scopelimit}[carried into every held-out result in \cref{sec:methods} and
\cref{sec:limitations}]
A held-out drug is not unseen by the whole pipeline. In natural-language
pretraining the language backbone may have encountered drug names and their
pharmacology, and every prompt supplies the drug name plus a \texttt{Mechanism}
field, an annotation where one exists and the literal string \texttt{unclear}
otherwise.
``Held out'' therefore means absent from Tahoe fine-tuning, not absent from the
model's prior world knowledge or from the inference prompt.
\end{scopelimit}

Holding data out comes at three levels of difficulty, not one. Trained
conditions test whether the model can reproduce information available during
fine-tuning. Unseen combinations hold out a whole treated well but often
recombine a drug and a cell line each encountered separately. Unseen drugs
remove all retained fine-tuning conditions for those drugs. These are not rungs
on one clean learning curve. They differ in support and in empirical
reliability, and only the last asks for a response without a Tahoe-trained
per-drug memory. The thesis therefore avoids calling all held-out examples
zero-shot and never pools the three strata into one transfer number.

% ---------------------------------------------------------------------------
\section{Four questions, each asked because of the answer to the last}
\label{sec:intro-map}

This thesis extends the cell-sentence line of work to chemical perturbation on
Tahoe-100M and asks whether a model fine-tuned on this atlas learns what a
particular drug does to a particular cell. A negative answer earns its place
only if it can be made specific, so the investigation asks four questions in
order, each one raised by the answer to the one before it.

\begin{enumerate}
  \item \textbf{Is the ruler trustworthy?} Can a predictor carrying no drug
        information score well on the field's standard metrics, and which
        metrics survive that control?
  \item \textbf{Does the model read the drug?} Measured on whichever metric
        survives the first question, with the input cell fixed and only the
        prompt's drug changed, does the prediction move at all?
  \item \textbf{Why does it behave that way?} Is the drug absent from the data,
        absent from the model's internal representation, or present in the
        representation but not consulted when the model generates?
  \item \textbf{Does redefining the effect repair it?} If the target is
        re-encoded as the drug-specific residual instead of the treated cell's
        own sentence, does the model become sensitive to the drug in its prompt?
\end{enumerate}

\noindent
The order is not presentational. A metric a zero-information predictor can
saturate cannot grade anything asked after it, and four of the five metrics
audited fail the corresponding control: the field's standard score is topped
outright by a predictor carrying no drug information at all, and in both audited
arms the other three rank a drug-agnostic baseline above a real positive
control. The survivor is a retrieval
test, defined next, and it is the repaired ruler every later answer is read
from. The null the second question returns on it is worth little until the third
locates it: drug identity is written into the activations and barely consulted
when the model generates. That diagnosis is what motivates the fourth, which
changes neither the model nor the data but the definition of what the model is
asked to produce.

% ---------------------------------------------------------------------------
\section{The question is identification, not resemblance}
\label{sec:intro-question}

Distributional similarity and token likelihood reward a prediction for looking
like a treated cell, even when the condition identity is wrong. The test left
standing is harder, and it is a retrieval test: among the other drugs measured
in the same context, does a prediction identify its own treatment?

\begin{defn}{Normalised inverse rank ($\NIR$)}
A prediction is compared with its own condition's truth and with the truths of
the other conditions in a declared comparison set. $\NIR$ is the fraction of that
set whose truth the prediction resembles \emph{less} than it resembles its own.
The comparison set is what fixes chance at $0.50$, in both measurement frames and
under an exchangeability condition that \cref{sec:res-metrics} audits with
empirical control arms. The score is tie-aware, and any component shared by every
candidate cancels.
\end{defn}

The metric is not infallible: its value depends on the gallery a prediction is
ranked against, it can compress near one, and its uncertainty must respect the
cell lines, drugs and wells that repeat in the design. What it will not do is
pass a prediction for resembling a generic treated profile, because every
treated profile does that.

% ---------------------------------------------------------------------------
\section{The bar is a dictionary, not chance}
\label{sec:intro-bar}

Two questions are routinely conflated in perturbation prediction. One is whether
a drug has a characteristic response that transfers across cellular contexts.
The other is whether a model uses the untreated cell to predict how \emph{this}
context departs from that average. Beating a control-copy told nothing about the
drug addresses neither, and approaching a split-half reference does not identify
which was solved.

The first has a model-free answer. A per-drug lookup predicts a
drug's residual as its mean residual over training conditions in other cell
lines. It needs no checkpoint and no prompt, and it measures the shared
component directly.\framecue{\Rframe}{scores here are on the drug-specific
residual, \cref{sec:methods-residual}} What it predicts is no small remainder.
Decomposing the treated-minus-control shift assigns $62.1\%$ of its squared
norm to the drug-specific residual and $39.0\%$ to the generic programme, with
a cross term of $-0.011$ closing the arithmetic.

The second is what a lookup cannot answer, because a per-drug constant makes the
same prediction in every context. How much of a drug's residual is shared across
contexts is measured in \cref{sec:res-transfer} by a disattenuated transfer
coefficient, \ci{\Tcoef=0.553}{0.514}{0.593}. The complementary
\ci{44.7\%}{40.7\%}{48.6\%} of that similarity scale is not shared across the
compared contexts. It is not a drug-by-cell-line interaction and not a variance
share; the estimator mixes cell line, dose, well and target-estimation
differences at once. It bounds what a per-drug constant leaves unresolved
without establishing that a model could recover any of it.

The win condition, fixed in advance, is therefore demanding. A
conditional model must beat the training-only drug lookup on common support;
only then would its access to the control cell have supplied information the
dictionary lacked. Beating chance would not do it, and neither would beating a
drug-agnostic control.

It was not met. On the $1{,}192$ conditions where model and lookup
are both defined, the lookup scores $0.913$ against the model's $0.549$. Its
advantage is not merely that it averages many cell lines; a lookup restricted to
one other cell line still scores $0.822$.\seealso{\cref{sec:res-lookup}: the
full baseline ladder, scored through one identical pipeline} Re-encoding the
target does make the model sensitive to the drug named in its prompt, which is
the project's first measurable trained-drug effect. Responding to the drug is
not the same as predicting what it does: on held-out treated wells the prompt
sensitivity survives, and biological transfer is \notestablished. Nothing built
here reaches the dictionary.

\begin{scopelimit}[carried into \cref{sec:res-metrics}, \cref{sec:res-lookup} and
\cref{sec:limitations}]
The within-well reference is a measuring instrument, not an oracle ceiling. It
asks how well one half of a treated population retrieves the other among
competing drugs. A model or a lookup may exceed it because it averages more
cells, borrows information across lines, or is compared with a noisy half-sample;
$\NIR$ also compresses close to one. The reference establishes empirical
discriminability at a stated cell budget and comparison set. It does not state
the maximum predictable fraction of biology, and it is recomputed within each
stratum and never transported across them.
\end{scopelimit}

Every conclusion that follows is re-tested under four changes to the
measurement: comparison sets are plate-matched, uncertainty is re-clustered on
drug, token budgets are equalised, and generic components are fitted without the
holdout they later grade. Not every conclusion survives them. Mechanism
annotation, for one, looks like a usable route to an unseen drug while
uncertainty is clustered on cell line, and stops looking like one once it is
re-clustered on the drug. A larger backbone was not tested, and a replication
with one is informative only if it receives the same split, targets, sentinels
and token budgets; is scored against the same
within-plate galleries and training-only lookups; and carries the same validity
and completion checks. Otherwise architecture, tokenisation, optimisation and
evaluation change together. This thesis does not supply that replication, but it
leaves an auditable ladder into which one can be inserted without weakening the
controls that produced the negative result.

\paragraph{The chapters} \Cref{sec:litreview} places these questions against
single-cell foundation models and the debate over whether deep perturbation
models beat simple baselines. \Cref{sec:methods-data} sets out the atlas, the
model and the rules of measurement; \cref{sec:methods} presents the investigation
as one chronological argument; \cref{sec:limitations} separates what is closed by
measurement from what remains untested; and \cref{sec:conclusions} concludes. The
symbols, terms and conventions are collected at the back in \nameref{sec:howtoread},
and nothing in the argument requires reading them first.
